\documentclass[11pt]{article}
\usepackage[hmargin=0.9in,vmargin=0.85in]{geometry}
\usepackage{amsmath,amssymb,amsthm}
\usepackage[T1]{fontenc}
\usepackage{lmodern}
\usepackage[hidelinks]{hyperref}
\setlength{\emergencystretch}{1em}
\newtheorem{theorem}{Theorem}
\newtheorem{lemma}{Lemma}
\newcommand{\F}{\mathbb F}
\newcommand{\N}{\mathbb N}
\newcommand{\Bl}{\operatorname{Bl}}
\title{One Modular Block Refutes the Half-Count Law\newline
  for Involution Centralizers}
\author{idealistichacker\\\small AI-assisted disproof of conjecture 00000005838}
\date{October 26, 2026}
\begin{document}
\maketitle
\begin{abstract}
Both source languages assert that an involution centralizer has half of its
blocks real-valued. Take the genuine involution in $G=C_2$ and the prime
$p=2$. Over $K=\overline{\F_2}$ the actual modular group algebra
$K[C_G(s)]$ has only the idempotents zero and one, and therefore precisely
one primitive central idempotent, or one modular block. Any selected
subfamily of this singleton has an integer cardinality, which cannot equal
$1/2$. We formalize this stronger predicate-parametric contradiction:
no definition of a predicate selecting blocks can make the asserted half
count true for this example. No ordinary-character classification or
invented realness test is required. All group, centralizer, field, algebra,
idempotence, primitivity, cardinality and half-count bridges are explicit.
\end{abstract}

\section{Exact Necessary Clause and Standard Block Convention}
The English source includes the clause ``the number of real-valued blocks
is half the total''; the Chinese source asserts the same half-count
equality. It is a separate necessary conjunct, not a conditional on an
unstated hypothesis about the surrounding dichotomy. Neither language
requires odd $p$, a nonabelian group or a simple group. We refute this
half-count clause and do not formalize or repair the source's undefined
kernel/dichotomy language, or separately refute its conjugate-pair clause.

The fixed complete bilingual source is retained separately:
\begin{quote}\small\ttfamily
Commit: b355e42906e6eaf52bc7e9d005e0ef02d8aa7400\\
Blob: bcf128f6c8b7e7c438e4e096e9e99e38c3135082.
\end{quote}
For a finite group $H$ and an algebraically closed field $K$ of
characteristic $p$, we use the standard modular algebra definition of its
$p$-blocks: nonzero primitive central idempotents of $K[H]$. Thus an
idempotent $e$ is central, satisfies $e^2=e$, and is primitive when it
cannot be written as a sum of two nonzero orthogonal central idempotents.
This is the standard block decomposition convention, with its connection
to the ordinary-character partition described in \cite[Chapter~12,
Sections~12.1--12.2]{Webb}. We do not instead count central idempotents of
the characteristic-zero semisimple algebra $\mathbb C[H]$.

The Lean block type is \emph{defined} by these actual ring equations and
the indecomposability condition. The general correspondence to an ordinary
character-partition API is cited and reviewed, not advertised as a Lean
character-theory theorem. No such theorem or realness classification is
needed for a contradiction applying to every predicate on the actual
modular block type. This standard algebraic definition, and not an
arbitrary Boolean labeling, fixes what is counted below.

\section{An Actual Involution Centralizer and Splitting Field}
Let $G=C_2=\{1,s\}$, with $s\ne1$ and $s^2=1$. It is a finite group
of order two. Since it is commutative,
\[
 C_G(s)=G.
\]
These are genuine group and subgroup statements, not an empty involution
case. In Lean, $G$ is \texttt{Multiplicative (ZMod 2)} and the centralizer
is the actual Mathlib subgroup of elements commuting with $s$.
The explicit group isomorphism $C_G(s)\simeq G$ induces an actual algebra
isomorphism $K[C_G(s)]\simeq K[G]$ by reindexing the group basis.
Idempotence and zero/one are transported across that isomorphism.

Take $p=2$ and the concrete field $K=\operatorname{AlgebraicClosure}(\F_2)$.
Its field, characteristic-two and algebraic-closedness instances are
constructed in Mathlib and explicitly audited here. The computation in
the next section works over \emph{every} characteristic-two field, so it
is valid over this actual splitting field, not just an enumeration over
the two-element field.

\section{All Idempotents of the Modular Group Algebra}
Write $[g]$ for the actual group-algebra basis vector at $g$. Every element
$x$ of $K[C_2]$ has a unique coefficient decomposition
\[
 x=a[1]+b[s],\qquad a,b\in K.
\]
This is proved for every group-algebra element using its actual finite
coefficient function and the two group elements. The multiplication rule
is $[g][h]=[gh]$, hence, in characteristic two,
\[
 x^2=(a^2+b^2)[1]+(ab+ba)[s]=(a^2+b^2)[1].
\]
If $x^2=x$, comparison of the $[s]$ coefficients first gives $b=0$.
Comparison at $[1]$ then gives $a^2=a$. Since $K$ is a field,
$a(a-1)=0$, so $a=0$ or $a=1$. Consequently
\[
 \forall x\in K[C_2],\quad x^2=x\ \Longrightarrow\ x=0\text{ or }x=1.
\]
This is a universal proof over the actual infinite splitting field, not
a conclusion from a finite grid or a stored count. Transport along the
centralizer algebra isomorphism gives the same result in $K[C_G(s)]$.

\begin{lemma}
If a nontrivial ring has only the idempotents zero and one, then it has
exactly one nonzero primitive central idempotent.
\end{lemma}
\begin{proof}
One is nonzero, central and idempotent. Suppose it decomposes into
orthogonal central idempotents $f,g$. If either is zero, the decomposition
is trivial. Otherwise both must equal one by the hypothesis, contradicting
$fg=0$. Thus one is primitive. Every nonzero primitive central idempotent
is a nonzero idempotent and hence equals one. This proves an actual
singleton block type, including existence, uniqueness and finiteness.
\end{proof}

Applying the lemma to $A=K[C_G(s)]$ yields
\[
 |\Bl(A)|=1.
\]
The ring is nontrivial because the coefficient field is nontrivial and
the group basis is nonempty. Lean proves finiteness of the full block
type and each predicate subtype; no infinite-type default cardinality
is used to obtain the count.

\section{No Selection of Blocks Can Satisfy the Half Count}
\begin{theorem}
For the above genuine involution centralizer, no predicate $P$ on its
modular blocks satisfies
\[
 |\{b\in\Bl(A):P(b)\}|=\frac{|\Bl(A)|}{2}.
\]
This holds with the equality interpreted in either $\mathbb Q$ or
$\mathbb R$.
\end{theorem}
\begin{proof}
The selected subtype has a natural cardinality $r$. Since the full block
count is one, the equality would imply $2r=1$, impossible for any
$r\in\N$. No value of $P$ has been assumed, chosen or hard-coded.
\end{proof}

In particular this applies to the actual real-block predicate, whether
real blocks are described as conjugation-stable blocks or via the
corresponding character family. It is unnecessary to decide that
predicate or classify the ordinary characters: the theorem excludes
\emph{every} possible block-selection predicate at once. The interpretation
of half is the ordinary rational/real equality, not truncated natural
division by two.

The formal final theorem allows an arbitrary selection family assigning a
predicate to the actual blocks of the centralizer of each $t\in C_2$.
It negates the half-count clause quantified over all nonidentity
involutions, by evaluating that clause at the displayed $s$. Both its
nonidentity premise and square-one premise are proved. Thus no arbitrary
choice of realness can make this necessary clause valid for this finite
group, and the source's universal conjunction is false. This argument
does not claim that other source clauses are separately false.

\section{Formal Scope and Reproducibility}
The four mathematical modules are BlockTheory, Involution, C2Algebra and
Counterexample. Their 34 public theorems include all coefficient,
idempotence, field, involution, centralizer/algebra transport, finite
cardinality and natural/rational/real half-count bridges.
The final declaration is
\texttt{Counterexample5838.half\_clause\_false\_for\_every\_selection}.
Its only explicit parameter is the arbitrary selection family; the
conclusion is its negation of the actual necessary clause. No desired
block count, realness value or negative conclusion is assumed.

Lean is pinned to 4.33.1 and Mathlib to
\begin{quote}\small\ttfamily
0df444a360eaa60ab8c11dca51a86af692955474.
\end{quote}
Every public theorem's compiled type and axiom footprint is audited.
Only the standard foundational axioms \texttt{propext},
\texttt{Classical.choice} and \texttt{Quot.sound} occur.
Finite exact group-algebra checks are regression evidence, not the
universal proof over $K$. Installed Lean and pinned dependency build caches
remain explicit trust inputs; no full-dependency clean build,
empty-machine bootstrap, cross-platform run or hosted-CI result is implied.
Separate AI review is not human expert endorsement, upstream acceptance,
priority, ranking or an award.

\begin{thebibliography}{9}
\bibitem{Webb} P. Webb, \emph{A Course in Finite Group Representation Theory},
  author-hosted prepublication text, Chapter~12, printed pp.~226--232,
  especially Section~12.1, Proposition~12.2.1 and Example~12.2.2.
\end{thebibliography}
\end{document}
